\documentclass[11pt]{article}
\usepackage[margin=1in]{geometry}
\usepackage[T1]{fontenc}
\usepackage{lmodern}
\usepackage{amsmath,amssymb,amsthm,mathtools}
\usepackage{booktabs,array,longtable}
\usepackage{enumitem}
\usepackage{microtype}
\usepackage{hyperref}
\usepackage{xcolor}
\setlength{\emergencystretch}{1.5em}
\usepackage{tikz}
\usetikzlibrary{arrows.meta,positioning,calc}

\hypersetup{hidelinks,pdftitle={Successive Concept Differentiation as Observation Refinement},pdfauthor={Brian Theory}}

\newtheorem{theorem}{Theorem}[section]
\newtheorem{proposition}[theorem]{Proposition}
\newtheorem{lemma}[theorem]{Lemma}
\newtheorem{corollary}[theorem]{Corollary}
\newtheorem{definition}[theorem]{Definition}
\newtheorem{example}[theorem]{Example}
\newtheorem{remark}[theorem]{Remark}
\newtheorem{question}[theorem]{Question}

\newcommand{\Ftwo}{\mathbb F_2}
\newcommand{\Cone}{\operatorname{Cone}}
\newcommand{\Nrv}{\operatorname{Nrv}}
\newcommand{\sk}{\operatorname{sk}}
\newcommand{\wt}{\operatorname{wt}}
\newcommand{\colim}{\operatorname*{colim}}
\newcommand{\im}{\operatorname{im}}
\newcommand{\coker}{\operatorname{coker}}
\newcommand{\Updown}{\mathbin{\Updownarrow}}

\title{\textbf{Successive Concept Differentiation as Observation Refinement}\\
\large Carrier Splitting, Interaction-Core Homology, and Descent}
\author{Brian Theory}
\date{Referee-repair draft v0.5 --- 29 September 2026}

\begin{document}
\maketitle

\begin{abstract}
A 2015 proposal for concept unfolding combined Boolean distinction, simplicial
boundary, and local-to-global reconstruction.  We separate these operations
and study one precise refinement model: a new Boolean observation splits the
states of a finite observational poset $Q$, producing an induced subposet
$Q^\psi\subseteq Q\times\{0<1\}$ with projection $\pi$.

Carrier splitting factors through a homotopy-neutral split completion followed
by order thinning.  In weighted height at most two its cone has a complete
integral rooted-incidence classification.  At arbitrary height, a two-colour
missing-copy model identifies the obstruction.  If the exclusive inversion
graph is a matching, cone homology is the direct sum of coarse edge-link
reduced homology shifted by two, with no torsion-freeness assumption.

Beyond matching, a filtration by mixed-core size has first page given by
available-link homology.  When cores have at most three vertices, an explicit
chain mapping cone computes the defect and the spectral sequence collapses at
$E^2$, with integral extensions retained.  Maximum-degree-two exclusive graphs
stabilize at $E^3$; an eight-state $K_{2,2}$ example proves this bound sharp.
A seven-state star rules out unconditional $E^2$ collapse for forests, Ferrers
graphs, and chordal-bipartite graphs.  An all-exclusive star family further
shows that interaction can create torsion even when all local links have
torsion-free homology.

The Boolean historical reconstruction and cover/descent framework are retained
as separate interpretations.  The general cofiber, obstruction-complex,
spectral-sequence, nerve, and sheaf machinery is classical.  The contribution
claimed here is the explicit observation-product specialization and its
positive classifications and integral failure boundaries; publication priority
for the exact specialization remains subject to review.

\end{abstract}

\tableofcontents

\section{Purpose and scope}

The motivating question is
\begin{center}
\fbox{\parbox{0.88\textwidth}{\centering Can successive differentiation of a concept be represented rigorously as a transformation between complexes?}}
\end{center}
The answer developed here is yes, provided that four levels are not conflated:
\begin{enumerate}[label=(\roman*)]
\item a semantic state space and its attributes;
\item an observation or distinction process;
\item a combinatorial/topological representation of the observed structure;
\item a genuine chain differential on that representation.
\end{enumerate}
The historical thesis is used only as a source of ideas.  The mathematical claims below stand on newly stated definitions.

The resulting architecture has two primary operators:
\[
 d:C_n(K_k;R)\longrightarrow C_{n-1}(K_k;R),\qquad d^2=0,
\]
and
\[
 U_k:C_\bullet(K_k;R)\longrightarrow C_\bullet(K_{k+1};R),
\]
where $d$ is the ordinary chain differential and $U_k$ is induced by an unfolding/refinement map.  When $U_k$ is a chain map, it induces homology maps; when the represented support changes, $H_*(\Cone(U_k))$ measures the change relative to the chosen representation.

\subsection{Conventions and result hierarchy}
Finite observational posets are assumed nonempty.  An order complex includes
the empty face.  Its augmented chain complex is denoted $\widetilde C_*$;
in particular $\widetilde H_{-1}(\Delta\varnothing;R)=R$.
Joins and suspensions of empty links use this simplicial convention; the
formal $(-1)$-sphere has no vertices.  Ordinary homology is zero in negative
degrees.  A coefficient ring is commutative with identity.

For a chain map $f:A_*\to B_*$, use
$\Cone(f)_n=B_n\oplus A_{n-1}$ and
$d(b,a)=(d_Bb+f(a),-d_Aa)$.
For a map of spaces or simplicial complexes, $\Cone(f)$ denotes its homotopy
cofiber and its homology is explicitly reduced.  For nonempty source and
target these conventions agree under the usual chain comparison.  A simplicial
map that repeats an image vertex sends that oriented simplex to zero on
normalized chains.  A shift $C[k]_n=C_{n-k}$ has differential $(-1)^k d_C$.

The main mathematical results are organized as follows.
\begin{enumerate}[label=(\arabic*)]
\item Theorem~\ref{thm:split-factorization} converts a changing observational
carrier into a neutral split completion followed by order thinning.
\item Theorem~\ref{thm:carrier-graph} gives the exact integral low-height result;
Theorem~\ref{thm:exclusive-matching} gives the all-height coarse-edge-link formula.
\item Theorems~\ref{thm:two-column-interaction} and
\ref{thm:degree-two-interaction} describe controlled overlap, including the
sharp distinction between $E^2$ and $E^3$ collapse.
\item Theorem~\ref{thm:main} treats cover refinement and descent separately.
It reconstructs data on the support actually covered; restriction from $X$ to
a proper subspace $Y$ need not preserve or uniquely extend global sections.
\end{enumerate}
The historical Boolean and simplex-skeleton constructions motivate the model;
they are not a claim to new general Boolean, polyhedral, or sheaf theory.

\section{Historical operator: what survives and what does not}

\subsection{Fixed-factor unfolding}

Let $A=\{a_1,\dots,a_m\}$ be a finite attribute set and let $\mathsf B(A)$ be the Boolean functions on $\{0,1\}^A$.

\begin{definition}[Factorized concept]
A factorized concept is an ordered tuple $F=(f_1,\dots,f_m)$ with $f_i\in\mathsf B(A)$.  Its conjunction semantics is
\[
\Pi(F)=\bigwedge_{i=1}^m f_i.
\]
\end{definition}

\begin{definition}[Historical one-step unfolding]
For a chosen factorization $F=(f_1,\dots,f_m)$, define
\[
U_{\mathrm{th}}(F)
=\bigvee_{j=1}^m\left(\neg f_j\wedge\bigwedge_{i\neq j}f_i\right).
\]
\end{definition}

\begin{proposition}[Exact-one-false characterization]
At a valuation $x$, $U_{\mathrm{th}}(F)(x)=1$ if and only if exactly one factor value $f_i(x)$ is false.
\end{proposition}
\begin{proof}
The $j$th disjunct requires $f_j(x)=0$ and every other factor to be $1$.  Distinct disjuncts are mutually exclusive.  Hence the disjunction is true exactly on the one-false layer.
\end{proof}

For the positive conjunction $X_1\wedge\cdots\wedge X_m$, the one-step output is therefore the Hamming layer with exactly one zero.  This is a precise mathematical version of ``one further distinction.''

\subsection{Representation dependence}

\begin{theorem}[Non-congruence of factorized unfolding]\label{thm:noncongruence}
Fix the arity $m=|A|\ge2$.  There is no Boolean-function operator
$\overline U:\mathsf B(A)\to\mathsf B(A)$ satisfying
\[
\overline U(\Pi(F))=U_{\mathrm{th}}(F)
\]
for every $m$-factor tuple $F\in\mathsf B(A)^m$.
\end{theorem}
\begin{proof}
The constant Boolean functions belong to $\mathsf B(A)$.  At the fixed
arity $m$, take
\[
 F=(0,1,1,\ldots,1),\qquad F'=(0,0,1,\ldots,1).
\]
Both conjunctions are the zero function.  Exactly one factor is false in
$F$, so $U_{\mathrm{th}}(F)=1$; exactly two are false in $F'$, so
$U_{\mathrm{th}}(F')=0$.  One operator on the conjunction cannot give both.
\end{proof}

For $m=1$, by contrast, $\Pi((f))=f$ and
$U_{\mathrm{th}}((f))=\neg f$; ordinary negation is such an operator.
Thus the failure of semantic descent begins at arity two.  For $m\ge2$,
either internal factorization/provenance is part of the concept, a canonical
normal form must be chosen, or the historical rule must be replaced by a
semantic operator.

\subsection{Boolean derivative interpretation}
For the monomial $f(x)=\prod_i x_i$, the Boolean partial derivative
\[
\partial_j^Bf=f(x_j=0)\Updown f(x_j=1)
\]
is $\prod_{i\neq j}x_i$.  Therefore
\[
U_{\mathrm{th}}(f)
=\bigvee_j\left((1-x_j)\wedge\partial_j^Bf\right).
\]
The historical rule is thus a negation-gated Boolean sensitivity operator, not the simplicial boundary.

\section{What the attached ``Logical Complexes'' notes recover}

The attached notes contain a useful idea, but two algebraic levels must be separated.

\subsection{XOR is a difference vector, not a scalar distance}
For Boolean state vectors $x,y\in\Ftwo^m$,
\[
x\Updown y=x+y
\]
is the vector difference in characteristic two.  The scalar Hamming distance is
\[
d_H(x,y)=\wt(x\Updown y).
\]
Thus expressions such as ``XOR is the distance'' should be replaced by: XOR is the coordinatewise disagreement vector, whose Hamming weight is the distance.

\subsection{The three-variable tetrahedron is genuine}
Set
\[
G=(1,1,1),\quad D=(1,1,0),\quad E=(1,0,1),\quad F=(0,1,1).
\]
Then
\[
D\Updown G=(0,0,1),\quad E\Updown G=(0,1,0),\quad F\Updown G=(1,0,0).
\]
These three differences form a basis of $\Ftwo^3$; viewed as points in $\mathbb R^3$, the corresponding real differences are also linearly independent.  Hence $\{G,D,E,F\}$ is an affine $3$-simplex, a tetrahedron.  The four points are not redundant.

A low-arity coincidence in the notes is also correct:
\[
D\Updown E=F,
\]
and cyclically.  Here $\{0,D,E,F\}$ is the even-parity plane in $\Ftwo^3$.  This closure does not persist for general $m$: XOR of two one-flip children of $1^m$ has Hamming weight $2$, not $m-1$.

\subsection{Why the signs disappear over \texorpdfstring{$\Ftwo$}{F2}}
The notes also anticipate a valid characteristic-two phenomenon.  If $[v_0,\dots,v_n]$ is a formal simplex, then over $\Ftwo$
\[
d[v_0,\dots,v_n]
=\sum_i[v_0,\dots,\widehat v_i,\dots,v_n]
\]
contains no visible alternating signs because $-1=+1$.  Every codimension-two face appears twice in $d^2$, so $d^2=0$.

The important correction is that the addition above occurs in the \emph{chain group on formal simplices}.  It is not XOR of the bit-vector labels themselves.  Label-space XOR and chain addition happen to use the same field in this model, but they are distinct operations on distinct vector spaces.

\section{Exact unfolding layers and hypersimplices}

Let the fully undifferentiated positive state be $\mathbf 1=(1,\dots,1)\in\{0,1\}^m$.  For $S\subseteq[m]$, define
\[
x_S=\mathbf 1\Updown \chi_S,
\]
so precisely the coordinates in $S$ are switched from $1$ to $0$.

\begin{definition}[Exact distinction layer]
The $k$th exact unfolding layer is
\[
L_k^{(m)}=\{x_S:|S|=k\}.
\]
\end{definition}

\begin{theorem}[Boolean unfolding simplex and hypersimplex theorem]\label{thm:hypersimplex}
For $m\ge2$:
\begin{enumerate}[label=(\roman*)]
\item $|L_k^{(m)}|=\binom{m}{k}$ and every $x\in L_k^{(m)}$ has Hamming weight $m-k$;
\item $\{\mathbf1\}\cup L_1^{(m)}$ is an affine $m$-simplex;
\item the convex hull of the exact $k$th layer is the hypersimplex
\[
\operatorname{conv}(L_k^{(m)})
=\left\{x\in[0,1]^m:\sum_i x_i=m-k\right\};
\]
\item for $0<k<m$, this polytope has dimension $m-1$ and is a simplex exactly in the extremal cases $k=1$ or $k=m-1$.
\end{enumerate}
\end{theorem}
\begin{proof}
Part (i) follows by choosing the $k$ zero coordinates.  For (iii),
the displayed intersection of the cube with the integer-sum hyperplane is a
bounded polytope.  If two coordinates of a point lie strictly between $0$
and $1$, perturb one upward and the other downward by a sufficiently small
amount; the point is not a vertex.  Exactly one fractional coordinate is
impossible because the coordinate sum is an integer.  Hence every vertex
is a $0$--$1$ vector with exactly $k$ zero coordinates, proving the hull
equality.  For (ii), if $u_i=\mathbf1\Updown e_i$, then $u_i-\mathbf1=-e_i$ in $\mathbb R^m$, so the $m$ difference vectors are linearly independent.  Part (iv) is the standard dimension and vertex description of the hypersimplex~\cite{GelfandEtAlHypersimplex}; in the two extremal nontrivial weights its vertices are the standard simplex vertices up to complementation.
\end{proof}

For $m=4$ the exact layers have the particularly transparent geometry
\[
\text{point}\;\longrightarrow\;\text{tetrahedron}\;\longrightarrow\;\text{octahedron}\;\longrightarrow\;\text{tetrahedron}\;\longrightarrow\;\text{point},
\]
since $\Delta(4,2)$ is an octahedron.  Thus repeated unfolding does not naturally stay inside the category of simplices if one geometrizes each exact state layer by convex hull.

\section{A canonical cumulative unfolding complex}

There is, however, a canonical simplicial model of \emph{distinction histories}.  A subset $S\subseteq[m]$ records the attributes that have been distinguished so far.

\begin{definition}[Cumulative distinction complex]
For $0\le r\le m-1$, define
\[
K_r^{(m)}=\{S\subseteq[m]:1\le |S|\le r+1\}.
\]
Thus $K_r^{(m)}=\sk_r\Delta^{m-1}$, the $r$-skeleton of the full simplex on the attribute set.  We also retain the empty face implicitly.
\end{definition}

The inclusions
\[
K_0^{(m)}\hookrightarrow K_1^{(m)}\hookrightarrow\cdots\hookrightarrow K_{m-1}^{(m)}=\Delta^{m-1}
\]
form a genuine filtration.  Here increasing $r$ means allowing larger jointly distinguished attribute sets.

\begin{theorem}[Exact cone profile of cumulative unfolding]\label{thm:skeletoncone}
Let $R$ be any coefficient ring and let
\[
i_r:K_{r-1}^{(m)}\hookrightarrow K_r^{(m)},\qquad 1\le r\le m-1.
\]
Then
\[
H_q\bigl(\Cone(C_*(i_r;R))\bigr)
\cong H_q(K_r^{(m)},K_{r-1}^{(m)};R)
\cong
\begin{cases}
R^{\binom{m}{r+1}},&q=r,\\
0,&q\neq r.
\end{cases}
\]
\end{theorem}
\begin{proof}
The relative chain complex $C_*(K_r,K_{r-1};R)$ has no new cells below degree $r$ and has one degree-$r$ basis element for every $(r+1)$-subset of $[m]$.  There are no relative cells above degree $r$.  Hence every relative differential is zero and the displayed relative homology follows.  The mapping cone of an inclusion computes the relative homology of the pair~\cite{Hatcher2002}.
\end{proof}

\begin{corollary}[Absolute homology of intermediate distinction stages]
For $0\le r<m-1$ and a field $\Bbbk$,
\[
\widetilde H_q(K_r^{(m)};\Bbbk)=0\quad(q\neq r),
\qquad
\dim\widetilde H_r(K_r^{(m)};\Bbbk)=\binom{m-1}{r+1}.
\]
The final stage $K_{m-1}^{(m)}=\Delta^{m-1}$ is contractible.
\end{corollary}

This gives a useful correction to the older suggestion that cycles should represent symmetry.  In this exact model the transient cycles are forced combinatorially by missing higher-order fillers.  They measure incomplete compatibility, not group-theoretic symmetry.

\begin{example}[Four attributes]
For $m=4$ the cumulative stages have Betti profiles
\[
K_0:(\beta_0)=(4),\qquad
K_1:(1,3),\qquad
K_2:(1,0,1),\qquad
K_3:(1,0,0,0).
\]
The relative/cone ranks at successive steps are $6$ in degree $1$, $4$ in degree $2$, and $1$ in degree $3$, exactly the numbers of newly admitted edges, triangles, and tetrahedra.
\end{example}

\section{Observation systems and three kinds of refinement}

The word ``refinement'' is overloaded.  Three operations should be distinguished.

\begin{definition}[Signature refinement]
An observation map is $q_k:X\to O_k$.  We say $q_{k+1}$ refines $q_k$ when there is $r_k:O_{k+1}\to O_k$ with $q_k=r_k\circ q_{k+1}$.  Equivalently, every $q_{k+1}$-fiber lies in a $q_k$-fiber.
\end{definition}

For finite propositional observation families, adding observables refines the generated topology, can split observational equivalence classes, thins the specialization preorder, and can increase signature distance.  This is an observation-theoretic notion of differentiation.

A second regime occurs when a new observation is constant on every current $T_0$ equivalence class.  Then the carrier does not split; only order relations can be removed.  This becomes a same-carrier order-thinning problem, which is treated separately in companion work.  If the new observation separates previously identified worlds, the refined carrier gains points and a different construction is required.

A third notion, used in the main theorem below, is a \emph{cover refinement}.  It concerns local pieces representing a support, not merely enlargement of an observation vocabulary.

\begin{remark}[Presentation-dependent nerve warning]
For an indexed family of truth regions $\{\varsigma(\varphi_i)\}$, the observation-side nerve records joint satisfiability, but it is presentation-dependent.  Adding a tautological observation can cone the nerve without changing the generated observation topology.  Therefore arbitrary extension of an observation list should not be identified with a topologically meaningful nerve filtration.  The main theorem uses actual refinement of covers with explicit containment data.
\end{remark}

\section{Carrier splitting under a new Boolean observation}\label{sec:carrier-splitting}

We now treat the refinement regime left open in the previous draft: a new
observation is not constant on every current observational-equivalence class, so
one old $T_0$ state can split into two refined states.  This operation is not a
cover refinement and is not, on its face, a same-carrier order thinning.  It does,
however, admit a canonical factorization through one.

Let $W$ be the finite world set, let $\Phi$ be the current observation family,
and write
\[
 Q=o_\Phi(W)\subseteq\{0,1\}^{I}
\]
for the finite $T_0$ quotient ordered coordinatewise.  Add one Boolean
observation $\psi$.  The refined quotient is
\[
 Q^\psi=o_{\Phi\cup\{\psi\}}(W)
 \subseteq Q\times\{0,1\},
\]
with the product order.  Let
\[
 \pi:Q^\psi\longrightarrow Q,
 \qquad \pi(s,a)=s,
\]
be the coordinate-forgetting map.  For $a\in\{0,1\}$ put
\[
 A_a:=\{s\in Q:(s,a)\in Q^\psi\},
 \qquad C:=A_0\cap A_1.
\]
Thus $A_0\cup A_1=Q$, and $C$ is precisely the set of old observational
states that split.

\begin{definition}[Split completion]
On the same carrier as $Q^\psi$, define the split-completion order
$\widehat Q^\psi$ by
\[
 (s,a)\leq_{\widehat Q}(t,b)
 \quad\Longleftrightarrow\quad
 \bigl(s<t\text{ in }Q\bigr)
 \ \text{or}\ 
 \bigl(s=t\text{ and }a\leq b\bigr).
\]
Equivalently, $\widehat Q^\psi$ is the lexicographic sum over $Q$ of the
nonempty fiber chains
\[
 F_s:=\{a\in\{0,1\}:(s,a)\in Q^\psi\}.
\]
Let $b(s,a)=a$ be the new-bit label on this expanded carrier.
\end{definition}

The distinction between $Q^\psi$ and $\widehat Q^\psi$ is exact.  The split
completion remembers every old comparison between coarse states, independently
of the new bit.  The actual refined order keeps only those comparisons along
which the new bit is nondecreasing.

\begin{theorem}[Carrier-splitting factorization]\label{thm:split-factorization}
With the notation above:
\begin{enumerate}[label=(\roman*)]
\item the projection $p:\widehat Q^\psi\to Q$, $p(s,a)=s$, is a homotopy
  equivalence of finite posets; in fact $\widehat Q^\psi$ beat-reduces to a
  subposet isomorphic to $Q$;
\item the actual refined quotient is exactly the one-bit thinning of the split
  completion,
  \[
    Q^\psi=(\widehat Q^\psi)_b,
  \]
  meaning that a comparison $x<_{\widehat Q}y$ is retained precisely when
  $b(x)\leq b(y)$;
\item if $i:\Delta(Q^\psi)\hookrightarrow\Delta(\widehat Q^\psi)$ is the
  inclusion, then $\Delta(\pi)=\Delta(p)\circ i$, and there is a homotopy
  equivalence of homotopy cofibers
  \[
  \operatorname{Cof}(\Delta\pi)\simeq
  \operatorname{Cof}(i).
  \]
  Consequently, for every coefficient ring $R$,
  \[
  \widetilde H_n\!\left(\Cone(\Delta\pi);R\right)
  \cong
  H_n\!\left(\Delta(\widehat Q^\psi),\Delta(Q^\psi);R\right).
  \]
\end{enumerate}
\end{theorem}

\begin{proof}
For every split state $s\in C$, the lower copy $(s,0)$ is an up-beat point of
$\widehat Q^\psi$: its strict upper set has the minimum $(s,1)$, because
$(s,1)$ lies below every copy over every $t>s$.  Remove these lower split
copies one at a time.  What remains has one representative over each $s\in Q$
and is order-isomorphic to $Q$.  Beat-point removal is a strong deformation
retraction for finite $T_0$ spaces~\cite{Stong1966}, giving (i).

For distinct coarse states $s<t$, the split completion always contains
$(s,a)<(t,b)$ whenever those copies exist.  The product order on $Q^\psi$
retains exactly the cases $a\leq b$.  The same statement holds within a split
fiber, where $(s,0)<(s,1)$ is retained.  This proves (ii).  Part (iii) follows
from $\pi=p\circ i$ and homotopy invariance of the homotopy cofiber under
postcomposition by the homotopy equivalence $\Delta(p)$.  Since $i$ is a
simplicial inclusion, its homotopy cofiber computes the relative homology of
its simplicial pair.
\end{proof}

\begin{remark}[What is standard and what is specific]
Beat-point reduction is classical finite-space topology~\cite{Stong1966}, and homotopy cofiber
invariance is standard; Quillen-type poset equivalences even admit simple-homotopy refinements in the standard setting~\cite{Barmak2011}.  The project-specific content is the observation-theoretic
factorization: an actual split of $T_0$ observational states is converted canonically
into a homotopy-neutral lexicographic expansion followed by the precise one-bit
monotonicity thinning already studied on a fixed carrier.  Thus carrier splitting
is not a fourth unrelated operation; it is a neutral blow-up plus an order defect.
\end{remark}

\subsection{The exact inversion-chain complex}

Theorem~\ref{thm:split-factorization} gives a chain-level classification without
any nerve hypothesis.

\begin{theorem}[Inversion-chain model]\label{thm:inversion-chain}
For any coefficient ring $R$, the relative group
\[
 C_k\!\left(\Delta(\widehat Q^\psi),\Delta(Q^\psi);R\right)
\]
is free on the strict chains
\[
 x_0<_{\widehat Q}x_1<_{\widehat Q}\cdots<_{\widehat Q}x_k
\]
whose bit word
\[
 b(x_0)b(x_1)\cdots b(x_k)
\]
is not weakly increasing.  Equivalently, a relative generator is exactly a
split-completion chain containing an inversion $1\cdots0$.  Its relative boundary
is the ordinary alternating simplicial boundary with every monotone-bit face set
to zero.
\end{theorem}

\begin{proof}
The two order complexes have the same vertex set.  A split-completion chain is a
simplex of $\Delta(Q^\psi)$ exactly when every comparison in the chain survives
the one-bit thinning, which is equivalent to the bit word being weakly increasing.
Quotienting the simplicial chain groups therefore leaves exactly the nonmonotone
chains, and the relative differential is the induced simplicial differential.
\end{proof}

This formulation is useful computationally: the defect depends only on the old
signature poset and the three-valued fiber profile
\[
 F_s\in\bigl\{\{0\},\{1\},\{0,1\}\bigr\},
\]
not on the number of worlds inside each old equivalence class.

\subsection{Immediate homotopy-neutrality criteria}

The factorization also reveals when splitting is automatically conservative.

\begin{proposition}[Monotone-section criterion]\label{prop:monotone-section}
If $A_1$ is an upward-closed subset of $Q$, then $\pi:Q^\psi\to Q$ is a
homotopy equivalence.  Dually, if $A_0$ is downward closed, then $\pi$ is a
homotopy equivalence.
\end{proposition}

\begin{proof}
Assume $A_1$ is upward closed.  Define
\[
 r_+(s)=
 \begin{cases}
 (s,1),&s\in A_1,\\
 (s,0),&s\notin A_1.
 \end{cases}
\]
The upset condition makes $r_+$ order preserving, and $\pi r_+=\mathrm{id}_Q$.
Moreover every $(s,a)\in Q^\psi$ satisfies $(s,a)\leq r_+(s)$, so
$\mathrm{id}_{Q^\psi}\leq r_+\pi$.  The order-homotopy lemma gives
$r_+\pi\simeq\mathrm{id}$ and hence $\pi$ is a homotopy equivalence.  The
$A_0$ statement is dual, using the minimum available bit in each fiber.
\end{proof}

\begin{corollary}[Common-outcome anchor]
If every old observational state has at least one refined world with $\psi=1$
($A_1=Q$), or every old state has at least one with $\psi=0$ ($A_0=Q$), then
carrier splitting is homotopy-neutral.  In particular, if every old state genuinely
splits, so $A_0=A_1=Q$, then $Q^\psi=Q\times\{0<1\}$ and the projection is a
homotopy equivalence.
\end{corollary}

Thus the act of duplicating a state is not itself the source of topological defect.
The defect is caused by incompatible outcome availability across comparable coarse
states: a lower state that can realize $1$ together with an upper state that can
realize $0$ creates a potentially deleted comparison.

\subsection{Localization by lower fibers}

For $s\in Q$ define the lower fiber
\[
 L_s:=\pi^{-1}(Q_{\leq s}).
\]
If $s\in A_1$, then $(s,1)$ is a maximum of $L_s$, hence $L_s$ is
contractible.  Therefore all lower-fiber obstructions to Quillen-type equivalence
localize to the $0$-only states $A_0\setminus A_1$.  Dually, all upper-fiber
obstructions localize to the $1$-only states $A_1\setminus A_0$.

This observation does not replace the exact relative complex above.  It does,
however, connect the carrier-splitting problem to the classical poset-fiber
framework.  Under the connectivity hypotheses of the Bj\"orner--Wachs--Welker
fiber theorem~\cite{BjornerWachsWelker2005}, the source $\Delta(Q^\psi)$ admits a wedge decomposition in terms
of the noncontractible lower fibers and upper links.  The important point here is
that the split fibers themselves disappear from that list whenever their top copy
is present.

\subsection{Exact low-height classification}

The general inversion complex can have arbitrarily high dimension.  In the first
nontrivial dimensional regime it collapses to graph incidence and admits a closed
integral classification.

\begin{definition}[Weighted split height]
Let $C=A_0\cap A_1$ be the split set.  Define
\[
 h_\psi(Q)
 :=\max_{s_0<\cdots<s_r\text{ in }Q}
 \left(r+\#\{j:s_j\in C\}\right).
\]
\end{definition}

\begin{lemma}\label{lem:split-height}
$h_\psi(Q)=\dim\Delta(\widehat Q^\psi)$.
\end{lemma}
\begin{proof}
A strict coarse chain contributes one vertex from every fiber, and a split fiber can
contribute both copies consecutively.  Taking both copies at every split state on the
chain attains the displayed number, and no split-completion chain can contain more.
\end{proof}

Suppose $h_\psi(Q)\leq2$.  An \emph{inversion edge} is a deleted order-complex edge
of $\Delta(\widehat Q^\psi)$, equivalently a pair
\[
 (s,1)<_{\widehat Q}(t,0),\qquad s<t,
\]
with $s\in A_1$ and $t\in A_0$.  A deleted triangle has one of the four nonmonotone
bit patterns
\[
 010,\qquad100,\qquad101,\qquad110.
\]
It therefore contains either one or two inversion edges.  Introduce one formal root
for each connected component of $Q$ (equivalently of $\widehat Q^\psi$).  The
\emph{carrier-split defect graph} $\mathcal D_\psi(Q)$ has one vertex for every inversion edge;
a deleted triangle with two inversion edges joins their vertices, while a deleted
triangle with one inversion edge joins that vertex to the root of its component.
Parallel graph edges are allowed.  Here a rooted forest means a forest
with exactly one formal root in each graph component.  This incidence graph
$\mathcal D_\psi$ is different from the exclusive inversion graph introduced later.

\begin{theorem}[Carrier-splitting cone classification in weighted height two]
\label{thm:carrier-graph}
Assume $h_\psi(Q)\leq2$.  Let $c_0(\mathcal D_\psi)$ be the number of connected components
of $\mathcal D_\psi(Q)$ containing no formal root.  Then
\[
 \widetilde H_n\!\left(\Cone(\Delta\pi);\mathbb Z\right)
 \cong
 \begin{cases}
   \mathbb Z^{\beta_1(\mathcal D_\psi)},&n=2,\\
   \mathbb Z^{c_0(\mathcal D_\psi)},&n=1,\\
   0,&\text{otherwise}.
 \end{cases}
\]
The relative boundary matrix is totally unimodular.  Hence the same Betti ranks hold
over every field.  Moreover the following are equivalent:
\begin{enumerate}[label=(\roman*)]
\item the carrier-splitting cone is acyclic over $\mathbb Z$;
\item it is acyclic over every field;
\item $\mathcal D_\psi(Q)$ is a rooted forest;
\item $\Delta(\widehat Q^\psi)$ collapses onto $\Delta(Q^\psi)$ through deleted
edge--triangle pairs;
\item $\Delta\pi:\Delta(Q^\psi)\to\Delta(Q)$ is a homotopy equivalence.
\end{enumerate}
\end{theorem}

\begin{proof}
The relative complex has only $C_2\to C_1$, since all vertices are retained
and the split completion has dimension at most two.  A deleted triangle of
bit type $100$ or $110$ has two inversion edges with opposite boundary signs;
a triangle of type $010$ or $101$ has one.  Consequently its matrix is the
oriented incidence matrix of $\mathcal D_\psi$ with formal-root rows deleted.
Each graph component contains at most one root, because triangles cannot
connect different coarse components.  Relative graph homology therefore gives
$H_2=\mathbb Z^{\beta_1(\mathcal D_\psi)}$ and
$H_1=\mathbb Z^{c_0(\mathcal D_\psi)}$.  Incidence matrices and their row
submatrices are totally unimodular.

Vanishing of both groups is exactly the rooted-forest condition.  In a rooted
tree with an edge, choose a nonroot leaf.  Its inversion edge occurs in just
one remaining deleted triangle and in no retained triangle, so the pair is
an elementary simplicial collapse.  Iterating leaf removal deletes every
relative edge and triangle.  Conversely a relative collapse gives a homotopy
equivalence and hence vanishing relative homology.  Finally the split-completion
projection is a homotopy equivalence, so two-out-of-three gives the statement
about $\Delta\pi$.  The free groups displayed above also show the equivalence
with acyclicity over every field.
\end{proof}

\begin{remark}[Status of the low-height theorem]
The graph-incidence algebra, total unimodularity, rooted forests, and beat-point
technology are standard.  The same height-two one-bit result is developed in the companion
order-thinning manuscript~\cite{DroncheffThinning2026}; the proof above makes
the transfer used here self-contained.  The new point here is the
carrier-splitting reduction that makes that theorem applicable to a process which
changes the $T_0$ carrier.  Accordingly, this theorem should be positioned as an
exact transfer theorem, not as a new general theory of rooted forests.
\end{remark}

\subsection{Sharpness beyond weighted height two}

No coefficient-independent forest criterion can hold in unrestricted dimension.
The carrier-splitting family contains the no-split case $C=\varnothing$, which is
exactly arbitrary same-carrier one-bit thinning.  Companion results~\cite{DroncheffThinning2026} show that one-bit
thinning in height three can already have relative integral torsion, and in unrestricted
height can encode arbitrary finite-poset homotopy types inside a contractible coarse
cone.  If one insists that at least one old state genuinely split, a disjoint neutral
split component may be adjoined without changing positive-degree cone homology.
Thus the weighted-height-two hypothesis in Theorem~\ref{thm:carrier-graph} marks a
real boundary rather than a presentational convenience.

The appropriate all-height object is therefore the inversion-chain complex of
Theorem~\ref{thm:inversion-chain}, supplemented where useful by poset-fiber or
homotopy-colimit methods.  The next two sections give all-height results under
matching and interaction-core hypotheses, together with explicit obstructions
to stronger graph-class conclusions.


\section{All-height carrier splitting under noninteracting exclusive inversions}
\label{sec:exclusive-matching}

The low-height rooted-incidence theorem gives a complete answer when the
weighted height is at most two.  In higher dimensions a universal rooted-forest
criterion is impossible: the same-carrier boundary case already supports
integral torsion and arbitrary finite-poset homotopy types.  Nevertheless, an
all-height exact classification survives when the genuinely incompatible
carrier outcomes do not interact.

Let
\[
A_a=\{s\in Q:(s,a)\in Q^\psi\},\qquad a\in\{0,1\},
\]
and put
\[
P=A_1\setminus A_0,\qquad Z=A_0\setminus A_1,\qquad
C=A_0\cap A_1.
\]
Thus $P$ consists of the $1$-only coarse states, $Z$ of the $0$-only coarse
states, and $C$ of the states that genuinely split.

\begin{definition}[Exclusive inversion graph]
\label{def:exclusive-inversion-graph}
The \emph{exclusive inversion graph} is the bipartite graph
\[
G_\psi=(P\sqcup Z,E_\psi),\qquad
E_\psi=\{(p,z)\in P\times Z:p<z\text{ in }Q\}.
\]
An edge records a coarse comparison whose earlier endpoint admits only the new
outcome $1$ while its later endpoint admits only $0$.
\end{definition}

\subsection{The full-doubling model}

Let
\[
D=Q\times\{0<1\},\qquad K=\Delta D,\qquad L=\Delta Q^\psi.
\]
The actual refined carrier is the induced subposet of $D$ obtained by deleting
exactly
\[
M_0=\{(p,0):p\in P\},\qquad
M_1=\{(z,1):z\in Z\}.
\]
Let $K_P$ be obtained by deleting $M_0$, let $K_Z$ be obtained by deleting
$M_1$, and put
\[
U=K_P\cup K_Z,\qquad L=K_P\cap K_Z.
\]
The complex $K_P$ closes upward to the full top row, while $K_Z$ closes
downward to the full bottom row.  Hence both $K_P\hookrightarrow K$ and
$K_Z\hookrightarrow K$ are homotopy equivalences.  Since $H_*(K,K_P)=H_*(K,K_Z)=0$, the triple $(K,K_P,L)$,
simplicial excision for $U=K_P\cup K_Z$, and the triple $(K,U,K_Z)$ give
$H_n(K,L)\cong H_n(K_P,L)\cong H_n(U,K_Z)\cong H_{n+1}(K,U)$.
Thus
\begin{equation}
H_n(K,L;R)\cong H_{n+1}(K,U;R)
\label{eq:two-colour-shift}
\end{equation}
for every coefficient ring $R$.  Since $D\to Q$ is a homotopy equivalence,
$H_n(K,L;R)$ is the reduced homology of the homotopy cofiber of the refined
projection $\Delta\pi:\Delta Q^\psi\to\Delta Q$.

A simplex of $K\setminus U$ contains at least one missing lower copy $(p,0)$
and at least one missing upper copy $(z,1)$.  Because it is a product-order
chain, every such pair satisfies $p<z$.  Therefore the missing vertices of a
relative simplex project to a complete bipartite subgraph of $G_\psi$.  This
observation isolates the source of higher interaction: bicliques with more than
one edge couple otherwise local defects.

\begin{proposition}[Empty-exclusive-graph neutrality]
\label{prop:empty-exclusive}
If $E_\psi=\varnothing$, then $\Delta\pi$ is a homotopy equivalence.
\end{proposition}
\begin{proof}
Set $b(s)=1$ exactly when some $p\in P$ satisfies $p\le s$.
This isotone choice takes an available value at every state: it is $1$ on $P$
and cannot be $1$ on $Z$ without an exclusive edge.  The section
$r(s)=(s,b(s))$ obeys $\pi r=\mathrm{id}$.
On $Q^\psi$, the map $c(s,a)=(s,\max\{a,b(s)\})$ is available and isotone,
with $\mathrm{id}\le c\ge r\pi$.  Comparable poset maps are homotopic,
so $r$ is a homotopy inverse.
\end{proof}

\begin{theorem}[Exclusive-Inversion Matching Theorem]
\label{thm:exclusive-matching}
Suppose that the exclusive inversion graph $G_\psi$ is a matching.  Then, for
every coefficient ring $R$ and every integer $n$,
\begin{equation}
\widetilde H_n\!\left(\Cone(\Delta\pi);R\right)
\cong
\bigoplus_{(p,z)\in E_\psi}
\widetilde H_{n-2}\!\left(
\operatorname{lk}_{\Delta Q}[p,z];R
\right).
\label{eq:exclusive-matching}
\end{equation}
Consequently $\Delta\pi$ is an $R$-homology equivalence if and only if every
coarse edge link $\operatorname{lk}_{\Delta Q}[p,z]$ indexed by an exclusive
inversion is $R$-acyclic.
\end{theorem}

\begin{proof}
A relative simplex of $(K,U)$ contains at least one vertex of $M_0$ and one
vertex of $M_1$.  If it contained two missing lower vertices and one missing
upper vertex, the corresponding upper endpoint would have degree at least two
in $G_\psi$; the dual statement holds for two missing upper vertices.  The
matching hypothesis therefore forces every relative simplex to contain a unique
mixed pair
\[
e=((p,0)<(z,1)),\qquad (p,z)\in E_\psi.
\]
Deleting either endpoint of $e$ sends the simplex into $U$, while deleting any
other vertex preserves the same pair.  Reorienting each simplex with the two core vertices first gives an even
chain shift.  Hence the relative chain complex splits
as
\[
C_m(K,U;R)=\bigoplus_{e\in E_\psi}C_m^{(e)},
\qquad
C_m^{(e)}\cong
\widetilde C_{m-2}(\operatorname{lk}_K e;R).
\]
Thus
\[
H_m(K,U;R)\cong
\bigoplus_{e\in E_\psi}
\widetilde H_{m-2}(\operatorname{lk}_K e;R).
\]

It remains to identify the product-order link.  For
$e=((p,0)<(z,1))$, its link factors into the lower interval, the open middle
interval, and the upper interval.  The open middle interval is covered by two
contractible subposets: delete $(z,0)$ and close upward to the bit-$1$ row, or
delete $(p,1)$ and close downward to the bit-$0$ row.  Their intersection is
$(p,z)_Q\times\{0<1\}$, homotopy equivalent to $(p,z)_Q$.  Since $(p,1)$ and
$(z,0)$ are incomparable, every middle-chain lies in one of these two pieces;
therefore
\[
\Delta((p,0),(z,1))\simeq\Sigma\Delta(p,z)_Q.
\]
Joining with the lower and upper factors gives
\[
\operatorname{lk}_{\Delta(Q\times\mathbf2)}e
\simeq
\Sigma\operatorname{lk}_{\Delta Q}[p,z].
\]
Combining this suspension shift with~\eqref{eq:two-colour-shift} yields
\eqref{eq:exclusive-matching}.
\end{proof}

\begin{corollary}[Bounded-poset localization]
\label{cor:bounded-localization}
Assume that $Q$ has a global minimum $\hat0$ and maximum $\hat1$, and that
$G_\psi$ is a matching.  Then every exclusive edge except possibly
$(\hat0,\hat1)$ has contractible coarse edge link.  Hence only the global
endpoint inversion can contribute to the cone homology.
\end{corollary}

\begin{proof}
If $p\ne\hat0$, the lower factor $\Delta Q_{<p}$ has the minimum $\hat0$; if
$z\ne\hat1$, the upper factor $\Delta Q_{>z}$ has the maximum $\hat1$.  In
either case the edge link is a join with a nonempty contractible complex.
\end{proof}

\begin{corollary}[Distributive endpoint profile]
\label{cor:distributive-endpoint}
Let $Q=J(R)$ be the lattice of order ideals of a nonempty finite poset $R$.  Make the
bottom state $1$-only, the top state $0$-only, and every other state split.  If
$R$ is an $r$-element antichain, so that $Q=B_r$, then
\[
\widetilde H_r(\Cone(\Delta\pi);\mathbb Z)\cong\mathbb Z
\]
and all other reduced cone homology vanishes.  If $R$ is not an antichain,
the cone is integrally acyclic.
\end{corollary}

\begin{proof}
Corollary~\ref{cor:bounded-localization} reduces the answer to the reduced
homology of the proper part of $J(R)$, shifted by two.  By the classical crosscut theorem~\cite{Bjorner1995}, the atom crosscut is a
full simplex unless the atoms join to the top element.  In $J(R)$ the atoms
correspond to the minimal elements of $R$, and they join to the top exactly
when every element of $R$ is minimal, i.e. exactly when $R$ is an antichain.
In that case $J(R)=B_r$ and its proper part is homotopy equivalent to
$S^{r-2}$.
\end{proof}

\subsection{Sharpness: interaction begins at the first shared endpoint}

The matching condition cannot be weakened merely to ``$G_\psi$ is a
forest.''  For the three-element chain $0<1<2$ with availability profile
\[
0:\{1\},\qquad 1:\{0\},\qquad 2:\{0\},
\]
the exclusive graph is the two-edge star $K_{1,2}$.  Direct integral Smith
normal form gives
\[
\widetilde H_1(\Cone(\Delta\pi);\mathbb Z)\cong\mathbb Z,
\]
while the sum of the two individual coarse-edge-link contributions is zero.
The shared endpoint creates a three-vertex interaction core and a differential
between the two edge defects.

Nor does matching imply total unimodularity or torsion-freeness.  Given any
nonempty finite poset $R$, take
\[
Q=\hat0\oplus R\oplus\hat1
\]
and use the endpoint profile above.  There is exactly one exclusive inversion,
and Theorem~\ref{thm:exclusive-matching} gives
\[
\widetilde H_n(\Cone(\Delta\pi);\mathbb Z)
\cong
\widetilde H_{n-2}(\Delta R;\mathbb Z).
\]
Thus arbitrary finite-poset homology, including torsion, can occur even inside
the matching class.  In particular, total unimodularity must be controlled by
stronger conditions on the relevant links; global torsion-freeness alone is
not a TU criterion~\cite{DeyHiraniKrishnamoorthy2011}.

The matching hypothesis is sufficient, not necessary: the same proof applies
whenever there is no mixed three-vertex core.  Section~\ref{sec:interaction-cores}
makes this criterion and the nonmatching interaction maps explicit.

\begin{remark}[Priority boundary]
The deletion/cofiber technology used here is classical.  In particular,
vertex-cover obstruction complexes already provide general homotopical tools
for comparing a complex with unions of induced subcomplexes
\cite{ChacholskiEtAl2021}, while poset-fiber decompositions give broad results
under fiber-connectivity hypotheses \cite{BjornerWachsWelker2005}.  The
candidate contribution of Theorem~\ref{thm:exclusive-matching} is the
observation-theoretic two-colour recognition and the resulting exact coarse
edge-link formula.  No claim is made that simplicial vertex deletion or
homotopy cofibers themselves are new.
\end{remark}


\section{Interaction cores beyond matching}
\label{sec:interaction-cores}

Retain $K=\Delta(Q\times\mathbf 2)$, $L=\Delta Q^\psi$, the missing
sets $M_0=P\times\{0\}$ and $M_1=Z\times\{1\}$, and
$U=K_P\cup K_Z$.  Cone homology in this section is reduced homology
of the topological homotopy cofiber, equivalently homology of the algebraic
mapping cone.  The two-colour shift is
\[
 \widetilde H_n(\operatorname{Cone}(\Delta\pi);R)
 \cong H_{n+1}(K,U;R).
\]
For a mixed missing core $S$, put
$\Lambda_S=\operatorname{lk}_K(S)[V(L)]$.
All link chain complexes below are augmented, including degree $-1$.

\begin{proposition}[Interaction spectral sequence]
\label{prop:interaction-ss}
The finite-filtration construction~\cite{Weibel1994} gives a convergent spectral sequence
\[
 E^1_{p,q}=\bigoplus_{|S|=p}\widetilde H_q(\Lambda_S;R)
 \quad\Longrightarrow\quad H_{p+q}(K,U;R),
\]
with $d^r$ of bidegree $(-r,r-1)$.  The first differential is the
signed sum of the inclusion maps obtained by deleting one missing vertex,
omitting nonmixed faces.  If all mixed cores have size at most $M\ge2$,
then $E^{M-1}=E^\infty$.
\end{proposition}
\begin{proof}
Orient a relative simplex by listing its $p$ missing vertices first and
its available vertices second, with the corresponding shuffle sign.
Its boundary is
\[
 d(S\wedge T)=\sum_i(-1)^i(S\setminus\{s_i\})\wedge T
                 +(-1)^pS\wedge\partial T.
\]
Terms with nonmixed core vanish in the relative complex.  Filtering by
$p$ gives the displayed page and differential.  The filtration is finite and exhaustive.  The abutment is the associated
graded of homology, so integral extension data may remain.
Only columns $2,\ldots,M$ can occur, so every $d^r$ with $r\ge M-1$
has zero source or target.
\end{proof}

\begin{theorem}[Exact two-column interaction theorem]
\label{thm:two-column-interaction}
Suppose no mixed core has four or more vertices.  Define
\[
 A_\bullet=\bigoplus_{|S|=2}\widetilde C_\bullet(\Lambda_S;R),
 \qquad
 B_\bullet=\bigoplus_{|S|=3}\widetilde C_\bullet(\Lambda_S;R),
\]
and let $h:B_\bullet\to A_\bullet$ be the signed sum of core-deletion
inclusions.  Then
\[
 C_\bullet(K,U;R)\cong\operatorname{Cone}(h)[2],\qquad
 \widetilde H_n(\operatorname{Cone}(\Delta\pi);R)
 \cong H_{n-1}(\operatorname{Cone}(h)).
\]
The interaction spectral sequence collapses at $E^2$, and there is a
natural exact sequence
\[
 0\longrightarrow\operatorname{coker}H_{n-1}(h)
 \longrightarrow\widetilde H_n(\operatorname{Cone}(\Delta\pi);R)
 \longrightarrow\ker H_{n-2}(h)\longrightarrow0.
\]
In particular the cone is $R$-acyclic if and only if $h$ is a
quasi-isomorphism.
\end{theorem}
\begin{proof}
The degree-$m$ relative group is $A_{m-2}\oplus B_{m-3}$ and its
boundary is $(a,b)\mapsto(\partial a+hb,-\partial b)$.
This is the asserted shifted chain cone.  Apply the two-colour shift,
the mapping-cone long exact sequence, and
Proposition~\ref{prop:interaction-ss} with $M=3$.
\end{proof}

The exact sequence splits noncanonically over a field.  Over $\mathbb Z$
it splits when the right-hand kernel is free abelian; no splitting is
asserted in general.  Thus collapse alone does not replace the integral
answer by a direct sum of link-homology kernels and cokernels.  Signed
incidence coefficients on link chains likewise do not imply total
unimodularity of the induced maps on homology.

\begin{theorem}[Degree-two classification and its sharp page bound]
\label{thm:degree-two-interaction}
Suppose the exclusive graph has maximum degree at most two.
Every mixed core has size at most four, so $E^3=E^\infty$.
If no $C_4$ component has both its same-colour pairs comparable in $Q$,
Theorem~\ref{thm:two-column-interaction} applies and $E^2=E^\infty$.
In every cycle component of length at least six the cone contribution is
the coarse-edge-link direct sum of the matching theorem.
\end{theorem}
\begin{proof}
Transitivity gives $N(p')\subseteq N(p)$ when $p<p'$ in $P$, and
$N(z)\subseteq N(z')$ when $z<z'$ in $Z$.
A mixed core is a biclique with each colour class a chain.  Degree two
bounds both class sizes by two.  A four-core requires a whole $K_{2,2}$
component with both colour pairs comparable.  This proves the bounds.
On a longer cycle, two same-colour degree-two vertices have distinct
neighbourhoods of the same cardinality and cannot be comparable if they
share a neighbour.  Hence there are no triple cores.  No missing vertex
can then lie in a product-edge link, so the original matching proof
applies componentwise without change.
\end{proof}

More explicitly, a path with at least three edges supports triples only
on its first two or last two edges: a comparable common-neighbour pair
must include a leaf.  For a path of at least four edges these end triples
have disjoint edge supports, giving independent star blocks and interior
edge blocks.  A three-edge path has up to two triples sharing its middle
edge and is treated by one instance of $\operatorname{Cone}(h)$.
Every surviving relative face remains in the same exclusive graph
component, so the component decomposition is at chain level.  Links
must nevertheless be calculated in the whole coarse poset.

For a two-edge star with ordered leaves, $h$ is the signed pair of
inclusions from its triple link to its two edge links.  On the frozen
three-chain profile $(1,0,0)$, all three available links have reduced
$H_0=\mathbb Z$ and $h_0$ is the primitive signed diagonal
$\mathbb Z\to\mathbb Z^2$.  This gives cone $H_1=\mathbb Z$,
although both coarse edge links are contractible.

\subsection{Higher differentials are necessary}
For a $C_4$ component whose two pairs are comparable, restrict to its
relative chain summand and write the core order
$a<b<x<y$.  Let $A,B,V$ be the sums of edge, triple, and four-core link
chains respectively.  Then
\[
 C_m(K,U)=A_{m-2}\oplus B_{m-3}\oplus V_{m-4},\qquad
 d(a,b,c)=(\partial a+hb,-\partial b+kc,\partial c),
\]
where $h,k$ are the signed core-deletion inclusions and $hk=0$.
In triple order $(abx,aby,axy,bxy)$, the signs of $k$ are $(-,+,-,+)$,
while $h(abx)=bx-ax$, $h(aby)=by-ay$, $h(axy)=ax-ay$, and
$h(bxy)=bx-by$, each term meaning its corresponding inclusion.
The only possible higher differential is
\[
 d^2:E^2_{4,q}\longrightarrow E^2_{2,q+1}.
\]
If $\partial c=0$ and $\partial b=kc$, it is represented by $hb$.
Thus a differential lowering core size by at most one can still produce
a nonzero secondary differential.

An explicit eight-state example makes this bound sharp.  Take
$Q=A_2\oplus R\oplus A_2$, with bottom antichain $\{a,a'\}$,
top antichain $\{y,y'\}$, and
$R=\{v,b,w,x\}$ with relations $v<x$, $b<w$, $b<x$.
Make $a,b$ one-only and $x,y$ zero-only; split the other states.
The only nonzero first-page groups are
$E^1_{2,1}=\mathbb Z$ and $E^1_{4,0}=\mathbb Z$.
The four-core link consists of $b_1,x_0$.  The triple links $abx$ and
$bxy$ are cones with apex $y'_1$ and $a'_0$, respectively; the links $aby$
and $axy$ are the two paths displayed below.  The edge links $ax$, $by$,
and $bx$ are cones, with apex $y'_1$, $a'_0$, and either of those vertices.
The remaining edge link $ay$ has maximal simplices
\[
 [a_1,b_1,w_1],\quad[v_0,x_0,y_0],\quad
 [a_1,v_1],\quad[v_0,v_1],\quad[w_0,w_1],\quad[w_0,y_0].
\]
Collapsing its two triangle ears leaves a cycle.  Thus the stated first-page
groups follow directly from all nine available links.  The paths
\[
 x_0-v_0-v_1-a_1-b_1,\qquad
 x_0-y_0-w_0-w_1-b_1
\]
in the two relevant triple links fill the four-core difference.
Give each path orientation from $x_0$ to $b_1$.  In the $ay$ component
of the signed zigzag, the first path minus the second traverses the loop once;
the other edge components are cycles in contractible links.  Therefore
$d^2$ is a unit $\mathbb Z\to\mathbb Z$, and $E^3=0$.
The supplementary signed-chain certificate verifies these fillings and
the integral homology directly.

Nor do Ferrers or chordal-bipartite graphs force $E^2$ collapse.
The seven-state poset with covers
$0<1$, $0<2$, $1<3$, $2<3$, $2<5$, $3<4$, $3<6$, $5<6$,
and profile $(1,2,0,0,2,2,0)$ has exclusive graph $K_{1,3}$.
Its integral $E^1=E^2$ has $\mathbb Z^2$ in bidegree $(2,1)$
and $\mathbb Z$ in $(4,0)$; $d^2$ is a primitive injection.
Its only nonzero cone group is $H_2=\mathbb Z$.
Here profile value $2$ means split.  The four-core link is
$\{0_1,2_0\}\simeq S^0$.  The three triple links, indexed by their
upper missing states $023$, $026$, and $036$, contain respective paths
\[
 A:2_0-4_1-0_1,\qquad B:2_0-5_1-0_1,\qquad
 C:2_0-3_0-1_0-1_1-0_1.
\]
All three triple links are contractible.  The $02$ edge link is the square
on bipartition $\{0_1,2_0\}\sqcup\{4_1,5_1\}$; the $03$ edge link is a
cone with apex $4_1$.  The $06$ edge link has maximal faces
\[
 [0_1,1_1],\ [0_1,5_1],\ [1_0,1_1],\
 [1_0,3_0,6_0],\ [2_0,3_0,6_0],\
 [2_0,5_0,5_1],\ [2_0,5_0,6_0].
\]
Collapsing its triangle ears leaves the cycle
$0_1-1_1-1_0-3_0-2_0-5_1-0_1$.
Orient $A,B,C$ from $2_0$ to $0_1$.  With the missing-core deletion
signs, their horizontal boundary has components $B-A$ in the $02$ link,
$A-C$ in the $03$ link, and $C-B$ in the $06$ link.  The first traverses
the $02$ square once; hence $d^2:\mathbb Z\to\mathbb Z^2$ has a unit
coordinate and is primitive.  This gives a signed witness independent of
the certificate's rank calculation.
For the abutment, write $T=Q\setminus\{0\}$ and $C=\{1,4,5\}$.
The complex $\Delta T$ is contractible: its two maximal-element lower ideals
are cones whose intersection is the lower ideal of $3$.
The complex $\Delta C$ is an edge and an isolated point.
In a star profile $Q=\{p\}\oplus T$, with $p$ one-only, $C$ split, and
$T\setminus C$ zero-only, the refined complex is the union of
$\Delta((T\times\{0\})\cup(C\times\{1\}))$ and the cone
$\Delta((\{p\}\times\{1\})\cup(C\times\{1\}))$.
The first piece retracts to $\Delta T$, and the intersection is $\Delta C$.
The homotopy-pushout description consequently gives
$\widetilde H_n(\Cone(\Delta\pi))\cong H_{n-1}(\Delta T,\Delta C)$.
Here this is $\mathbb Z$ only at $n=2$; the sole possible $d^2$ must be a
primitive injection.  The supplementary certificate records all link matrices
and signed chains.  These are explicit counterexamples, not global minimality claims.

\subsection{An unbounded interacting class and new torsion}
\begin{proposition}[All-exclusive star]
Let $R$ be a nonempty finite poset.  Take $Q=\{p\}\oplus R$, make
$p$ one-only and every state of $R$ zero-only.  Then $E^2=E^\infty$
integrally and
\[
 \widetilde H_n(\operatorname{Cone}(\Delta\pi);\mathbb Z)
 \cong H_{n-1}(\Delta R;\mathbb Z),
\]
where the right side is ordinary homology.
\end{proposition}
\begin{proof}
Every core is $\{p_0\}\cup\sigma_1$ for a nonempty chain $\sigma$
in $R$.  Its available link is the isolated vertex $p_1$ disjoint from
the cone $\Delta(R_{\le\min\sigma}\times\{0\})$.
The difference-of-components generators identify all nonzero link
homology with $\mathbb Z$ in degree zero and identify their inclusion
maps with the identity.  Consequently $d^1$ is, up to a uniform sign,
the ordinary simplicial boundary of $\Delta R$.  The single row
collapses on $E^2$ and has only one filtration grade in each total degree.
\end{proof}

For the nonempty-face poset of a triangulated projective plane, all these
available links have torsion-free homology, and every coarse edge link
is a vertex link in its barycentric subdivision, hence a circle.
Nevertheless the interaction differential creates $\mathbb Z/2$ in
cone degree two.  The explicit six-vertex triangulation in the
supplement gives a 32-state example, with every link and the reduced
interaction complex checked by integral Smith normal form.

\subsection{Whole-relation Horn closure is neutral}
\begin{proposition}
If the nonempty refined Boolean relation is closed under coordinatewise meet,
then $\Delta\pi$ is a homotopy equivalence.  The dual statement holds for
coordinatewise join closure.
\end{proposition}
\begin{proof}
Finiteness and meet closure give a minimum in the refined relation.
Its coordinate projection is again meet closed and nonempty, so also has a
minimum.  Both order complexes are contractible.  Join closure gives maxima.
\end{proof}
The hypothesis concerns the whole realized refined relation; it is not merely
a closure assumption on a selected subset of observations.  This is an
elementary application of closure-system topology, not a new Horn theorem.

\begin{remark}[Scope and priority]
Filtered-complex spectral sequences, link/cofiber recursions, and their
extension issues are classical.  The complexes $\Lambda_S$ are precisely
vertex-cover obstruction complexes in the sense of
Chach\'olski--Jin--Scolamiero--Tombari~\cite{ChacholskiEtAl2021}.
The results here explicitly identify which such complexes interact in
the observation product model, with positive class restrictions and
integral obstructions.  They do not introduce a new general obstruction
theory or supersede poset-fiber theorems.
\end{remark}


\section{Cover refinement and canonical homology maps}

\begin{definition}[Nerve]
For a finite cover $\mathcal U=\{U_i\}_{i\in I}$ of $X$, its nerve $\Nrv(\mathcal U)$ has vertex set $I$ and simplex $\sigma\subseteq I$ exactly when
\[
U_\sigma:=\bigcap_{i\in\sigma}U_i\neq\varnothing.
\]
\end{definition}

\begin{definition}[Refinement map]
Let $\mathcal V=\{V_j\}_{j\in J}$ and $\mathcal U=\{U_i\}_{i\in I}$ cover the same space, or let the union of $\mathcal V$ lie in the union of $\mathcal U$.  A refinement choice is a map $\rho:J\to I$ such that
\[
V_j\subseteq U_{\rho(j)}
\]
for every $j$.
\end{definition}

\begin{lemma}[Contiguity of refinement choices]\label{lem:contiguity}
Every refinement choice induces a simplicial map
\[
N(\rho):\Nrv(\mathcal V)\to\Nrv(\mathcal U).
\]
Any two refinement choices $\rho,\rho'$ are contiguous and therefore induce the same homomorphism on homology.
\end{lemma}
\begin{proof}
If $\sigma\in\Nrv(\mathcal V)$, choose $x\in\bigcap_{j\in\sigma}V_j$.  Then $x$ belongs to every $U_{\rho(j)}$ and every $U_{\rho'(j)}$.  Hence the union of the two image vertex sets spans a simplex of $\Nrv(\mathcal U)$, which is precisely contiguity.
\end{proof}

This lemma is standard.  In particular, Leit\~ao's 2026 framework~\cite{Leitao2026} for persistent homology of cover refinements makes this contiguity principle explicit at the cover level.  It should not be presented here as a new theorem.

\section{Observation--Refinement Descent and Cone Theorem}

We now state the central structural result in a form adapted to successive concept differentiation.

\begin{definition}[Pair-good subcomplex cover]
Let $Y\subseteq X$ be a finite simplicial pair.  A finite family
$\mathcal U=\{U_i\}_{i\in I}$ of subcomplexes covering $X$ is
\emph{good relative to $Y$} if every nonempty intersection
$U_\sigma=\bigcap_{i\in\sigma}U_i$ is contractible and every nonempty
$U_\sigma\cap Y$ is contractible.  The restricted cover
$\mathcal U|_Y=\{U_i\cap Y\}_{i\in I}$ is indexed after deleting its
empty members.  A finite subcomplex cover of $Y$ is \emph{good} if all
its nonempty intersections are contractible.  The subcomplex condition
is part of these definitions, not a consequence of intersection
contractibility alone.
\end{definition}

\begin{theorem}[Observation--Refinement Descent and Cone Theorem]\label{thm:main}
Let $Y\subseteq X$ be a finite simplicial pair.  Let $\mathcal U$
be a finite pair-good subcomplex cover of $X$ and let $\mathcal V$ be
a finite good subcomplex cover of $Y$ refining $\mathcal U|_Y$.
Delete empty indexed members before taking nerves.  Thus all three covers
$\mathcal U$, $\mathcal U|_Y$, and $\mathcal V$ lie in the
functorial subcomplex-cover nerve theorem.  Choose any refinement map
\[
\rho:\mathcal V\longrightarrow\mathcal U|_Y
\]
and let
\[
u:
\Nrv(\mathcal V)\xrightarrow{N(\rho)}\Nrv(\mathcal U|_Y)
\hookrightarrow\Nrv(\mathcal U)
\]
be the induced simplicial map.  Then:
\begin{enumerate}[label=(\alph*)]
\item the homotopy class up to contiguity, and hence the induced homology map $u_*$, is independent of the refinement choice;
\item $N(\rho):\Nrv(\mathcal V)\to\Nrv(\mathcal U|_Y)$ is a homotopy equivalence;
\item for any coefficient ring $R$,
\[
H_n\!\left(\Cone(C_*(u;R))\right)
\cong H_n(X,Y;R)
\]
naturally up to the standard nerve identifications;
\item in the conservative case $Y=X$, the cone is acyclic, equivalently $u$ is a quasi-isomorphism;
\item assume additionally that $X$, $Y$, the cover members, and their
intersections are admissible objects of a site with inclusion
morphisms for the displayed subcomplexes, their pullbacks are represented
by these intersections, and the displayed families are covers in that
site.  Take the sheaf and cosheaf values in a category admitting the
finite limits and colimits used below, for example abelian groups.
For a sheaf $\mathcal F$, compatible data
on $\mathcal U$ reconstruct $\mathcal F(X)$; compatible data on
$\mathcal U|_Y$ or $\mathcal V$ reconstruct $\mathcal F(Y)$.
Refinement between the latter two covers preserves that section on $Y$.
No unique extension from $Y$ to $X$ is asserted;
\item under the same site hypothesis, a cosheaf $\mathcal G$ is assembled
as the colimit of the intersection diagram on each covered support: the
colimits are $\mathcal G(X)$ and $\mathcal G(Y)$, respectively.
Refinement on $Y$ preserves its colimit; inclusion induces the comparison
$\mathcal G(Y)\to\mathcal G(X)$, which need not be an isomorphism.
\end{enumerate}
\end{theorem}

\begin{proof}
Part (a) is Lemma~\ref{lem:contiguity}.  For (b), the finite
subcomplex-cover nerve theorem~\cite{BauerEtAl2023,Ramras2026}
supplies compatible equivalences for $\mathcal V$ and
$\mathcal U|_Y$ with $Y$.  Functoriality for the refinement map
identifies $N(\rho)$ with the identity of $Y$ through these
equivalences.  Hence $N(\rho)$ is a homotopy equivalence.

Let $j:\Nrv(\mathcal U|_Y)\hookrightarrow\Nrv(\mathcal U)$.  Since $N(\rho)$ is a homotopy equivalence, the mapping cones of $j\circ N(\rho)$ and $j$ have isomorphic homology; this follows, for example, by comparing the long exact sequences of the two cones and applying the five lemma.  The mapping cone of the chain inclusion $C_*(j)$ computes
\[
H_n(\Nrv(\mathcal U),\Nrv(\mathcal U|_Y);R).
\]
The relative nerve theorem identifies this with $H_n(X,Y;R)$, proving (c).  If $Y=X$, the relative groups vanish, giving (d).

For (e), apply the sheaf axiom~\cite{Curry2013} separately on each support: compatible
families on a cover of $S$ are the equalizer describing $\mathcal F(S)$,
for $S=X$ or $S=Y$.  The covers $\mathcal U|_Y$ and $\mathcal V$ cover
the same $Y$, so their gluing maps agree with refinement.  Restriction
$\mathcal F(X)\to\mathcal F(Y)$ is a separate map and need not be bijective.
Part (f) is the dual cosheaf colimit statement, again on each support separately.
\end{proof}

\begin{example}[Restriction is not unique extension]
Take the discrete space $X=\{a,b\}$, $Y=\{a\}$, and the sheaf of
integer-valued functions.  Singleton covers satisfy the stated good-cover
conditions.  Data on $Y$ determine an element of $\mathbb Z$, whereas
$\mathcal F(X)=\mathbb Z^2$; restriction is the projection onto the first
coordinate.  A section on $Y$ therefore has many extensions to $X$.
This distinguishes descent on a fixed support from an extension problem.
\end{example}

\subsection{What the theorem means for unfolding}
The theorem separates two questions that were previously mixed:
\begin{enumerate}[label=(\roman*)]
\item \emph{Does the chosen map fail to be a homology equivalence?}  Cone/relative homology detects this failure; its vanishing alone does not detect every change of support or homotopy type.
\item \emph{Can compatible local descriptions be reconstructed globally?}  This is answered by descent or colimit conditions.
\end{enumerate}
A finer description can therefore be topologically conservative while still providing more local information.  Conversely, one may have perfect local gluing on a support whose topology has genuinely changed.

\section{Descent defects for data that are not already a sheaf}

If one assumes a sheaf from the beginning, reconstruction is built into the axioms.  For modeling it is useful to measure how a candidate local-data system fails to be a sheaf.

Let $\mathcal F$ be a presheaf of abelian groups and $\mathcal U=\{U_i\}$ a finite cover.  Define
\[
C^0(\mathcal U;\mathcal F)=\prod_i\mathcal F(U_i),\qquad
C^1(\mathcal U;\mathcal F)=\prod_{i<j}\mathcal F(U_i\cap U_j),
\]
and
\[
(\delta^0s)_{ij}=s_j|_{U_i\cap U_j}-s_i|_{U_i\cap U_j}.
\]
Let $Z^0=\ker\delta^0$ be the compatible local families and let
\[
\operatorname{res}:\mathcal F(X)\to Z^0(\mathcal U;\mathcal F)
\]
be global restriction.

\begin{definition}[Descent defects]
Define
\[
D_{\mathrm{uniq}}(\mathcal U;\mathcal F)=\ker(\operatorname{res}),
\qquad
D_{\mathrm{exist}}(\mathcal U;\mathcal F)=\coker(\operatorname{res}).
\]
\end{definition}

\begin{proposition}[Exact reconstruction criterion]
Compatible local data reconstruct a unique global element if and only if
\[
D_{\mathrm{uniq}}=0=D_{\mathrm{exist}}.
\]
In particular both defects vanish for a sheaf.
\end{proposition}

This motivates a two-axis defect signature for an unfolding step $u$:
\[
\mathfrak D(u;\mathcal F)
=\left(
H_*(\Cone C_*(u)),\,
D_{\mathrm{uniq}},\,
D_{\mathrm{exist}}
\right).
\]
The first entry is structural; the latter two are semantic/local-to-global.  They should not be collapsed into one number.

\section{Exact simplicial gluing and the no-filler warning}

\begin{proposition}[Exact gluing as a colimit]
If a simplicial complex $K$ is covered by subcomplexes $\{K_i\}_{i\in I}$, then $K$ is the colimit of the full diagram consisting of the $K_i$ and all their finite intersections, with the evident inclusion maps.
\end{proposition}

This is the rigorous combinatorial version of rebuilding an object from pieces with overlap provenance.

\begin{proposition}[Bare proper faces do not determine a filler]
The proper-face data of an $n$-simplex do not by themselves determine whether the $n$-simplex is present.
\end{proposition}
\begin{proof}
The boundary complex $\partial\Delta^n$ and the filled simplex $\Delta^n$ have exactly the same proper faces, but only the latter contains the top-dimensional simplex.
\end{proof}

Therefore the historical idea of an ``integrator'' requires a filler rule, provenance, a flag/clique convention, a colimit diagram containing the intended top cell, or semantic descent data.  Conjunction of face-formulas is not a general simplicial gluing operation.

\section{Finite prototypes}

\subsection{Three attributes: vertices, cycle, filler}
The cumulative distinction filtration for $m=3$ is
\[
K_0=\{1,2,3\}
\hookrightarrow
K_1=\partial\Delta^2
\hookrightarrow
K_2=\Delta^2.
\]
Over $\Ftwo$:
\[
H_0(K_0)\cong\Ftwo^3,
\qquad
H_1(K_1)\cong\Ftwo,
\qquad
H_{>0}(K_2)=0.
\]
The first step adds three edges and has
\[
H_1(K_1,K_0)\cong\Ftwo^3.
\]
The second adds one triangle and has
\[
H_2(K_2,K_1)\cong\Ftwo.
\]
The one-dimensional cycle at the middle stage is therefore not a symmetry invariant; it is the obstruction created by having all pairwise compatibilities without the triple filler.

\subsection{Four attributes: a tetrahedral history complex}
For $m=4$:
\[
K_0\subset K_1\subset K_2\subset K_3=\Delta^3,
\]
where $K_1$ is the complete graph $K_4$ and $K_2$ is the tetrahedral boundary.  The nonzero reduced homology is
\[
\widetilde H_0(K_0)\cong R^3,\qquad
H_1(K_1)\cong R^3,\qquad
H_2(K_2)\cong R.
\]
The next stage fills the tetrahedron and kills the $2$-cycle.

At the exact-state level, however, the $k=2$ Hamming layer has six points and convex hull an octahedron.  This illustrates why the state-layer geometry and the distinction-history complex are related but not identical constructions.

\section{Symmetry, indistinguishability, and homology are different invariants}

Let $\operatorname{Aut}(K)$ be the simplicial automorphism group.  Let an observation map $q:X\to O$ define indistinguishability by equal fibers.  Neither is determined by homology.

A filled triangle has automorphism group $S_3$ but trivial positive-dimensional homology.  Conversely, an asymmetric unicyclic graph can have $H_1\cong\mathbb Z$ and trivial automorphism group.  Thus cycles cannot generally be identified with symmetry.

There is nevertheless a rigorous symmetry-breaking statement at the observational level.  If a group $G$ acts on $X$, define
\[
G_q=\{g\in G:q(gx)=q(x)\ \forall x\in X\}.
\]
If $q_{k+1}$ refines $q_k$, then
\[
G_{q_{k+1}}\le G_{q_k}.
\]
Finer observation can therefore break observational symmetry by shrinking the subgroup that preserves all observations.  This is separate from the homology of the chosen complex.

\section{Relation to observation topologies and order thinning}

The observation-topology program provides a natural semantic base for the present construction.  An indexed family $\Phi=(\varphi_i)$ defines signatures and an observation-generated topology.  Extending $\Phi$ by new coordinates makes the topology finer, refines observational equivalence, and thins the specialization preorder.  These facts supply a mathematically clean interpretation of ``more differentiated observation.''

However, the passage from observation topology to complex must be chosen explicitly.  The nerve of the indexed truth regions records joint satisfiability but is not an invariant of the generated topology.  This is why the present paper uses a cover-refinement contract when applying the nerve theorem.

There is also an important same-carrier subcase.  If a new Boolean observation factors through the current $T_0$ quotient, it does not split observational states; it only removes comparisons incompatible with the new bit.  Companion work shows that, for a height-two current signature poset and one new bit, homotopy neutrality can be recognized exactly by a rooted-forest criterion and certified by simplicial collapse.  That result complements the present theorem: the present paper treats changing carriers, matching and interacting cone defects, and cover-level descent, while the companion order-thinning paper develops the same-carrier incidence theorem and its separate sharpness results.

\section{Persistence and zigzags}

A monotone sequence of inclusions of concept complexes
\[
K_0\hookrightarrow K_1\hookrightarrow\cdots
\]
produces an ordinary persistence module on homology.  If differentiation alternates between additions, deletions, support restrictions, or reinterpretations, the natural object is instead a zigzag
\[
K_0\to K_1\leftarrow K_2\to K_3\leftarrow\cdots.
\]
This is preferable to forcing every cognitive or observational change into a monotone filtration.

Cover-level persistence is already an active theory.  In particular, Leit\~ao develops persistent homology directly from cover refinements and proves stability statements that propagate through nerve and co-nerve functors.  Therefore any novelty claim here must be made at the level of the specific observation semantics, relative defect interpretation, or descent coupling, not at the level of cover-refinement persistence itself.

\section{Priority and novelty position}
The standard ingredients include Boolean/Hamming and hypersimplex geometry,
chain cones, beat-point reductions, relative nerve theorems, sheaf descent,
poset-fiber methods, and filtered-complex spectral sequences.
For the cover $X_0=V(K)\setminus M_0$, $X_1=V(K)\setminus M_1$, with
$A=X_0\cap X_1=V(L)$, each $\Lambda_S$ is precisely the available-vertex
obstruction complex $\operatorname{St}(S,A)$ of
Chach\'olski--Jin--Scolamiero--Tombari~\cite{ChacholskiEtAl2021}, with
our augmented empty-face convention.  Their published Corollary~7.2 treats
several outside vertices and Theorem~8.6 controls the homotopy fibers of the
vertex-cover inclusion through these obstruction complexes.

The Bj\"orner--Wachs--Welker fiber theorem~\cite{BjornerWachsWelker2005}
gives decompositions under explicit fiber-connectivity or null-attachment
hypotheses.  Here the attachment maps are retained and may produce torsion
or higher differentials.  This is a narrower exact computation, not a
replacement for general fiber theory.  Likewise, Ferrers graph resolutions,
graph independence complexes, and unweighted biclique complexes do not by
themselves identify the available-link diagram used here.

The candidate contribution is the explicit observation-product recognition,
the coarse-edge-link matching formula, the transitivity restrictions on
interacting cores, and the sharp integral positive and negative examples.
The matching deletion step is a direct specialization of the obstruction
decomposition above; identifying its product-edge link with a suspension
of the coarse link is the model-specific calculation.  The interaction
first page, filtration-width collapse, and two-column kernel/cokernel
sequence use standard filtered-chain and mapping-cone algebra after the
available-link diagram has been identified.  These tools are not claimed
as independent new constructions.
The historical audit and descent framework provide interpretation rather
than priority for classical machinery.  No claim of first discovery of the
exact specialization is made solely from its absence in a targeted search.
The separate convergence report records the primary-source comparison and
the remaining publication-level judgment.

\section{Next theorem target}
The matching case, the two-column interaction case, and the degree-two
$E^3$ bound are now established.  Unrestricted $E^2$ collapse in the
proposed graph classes is false, so it is no longer an open target.

\begin{question}[Integral secondary maps in a four-cycle]
For each integer $k\ge2$, can an exclusive $K_{2,2}$ profile with
torsion-free available-link homology realize
$d^2:\mathbb Z\to\mathbb Z$ as multiplication by $k$?
Alternatively, what restrictions does the product-poset geometry impose
on the possible integral secondary maps?
\end{question}
This isolates interaction-created torsion at the smallest exclusive graph
that supports a four-core.  The unit example above does not resolve it.
The smallest possible higher-differential or interaction-torsion realization
also remains open globally; induced-restriction minimality is weaker.

\section{Conclusion}
Observation refinement is a transformation between representations, not a
chain boundary.  Carrier splitting has a neutral-completion factorization,
an exact low-height incidence description, and an all-height matching
edge-link formula.  Beyond matching, available-link inclusions form the
interaction complex.  Two core columns give an exact chain cone and
$E^2$ collapse; degree two permits a genuine $d^2$ and only the sharp
$E^3$ bound in general.  Interaction can create integral torsion even from
torsion-free local links.

Cover refinement and descent answer a separate question on the support
actually covered.  The historical Boolean reconstruction, local-to-global
data, and structural homology therefore remain distinct.  The mathematical
claims of this draft are the stated specialized formulas and failure
boundaries; general topological machinery and publication priority are not
silently strengthened.

\appendix
\section{Reproducibility}
The accompanying script \texttt{verify\_unfolding\_complexes.py} checks the following finite claims over $\Ftwo$ for $2\le m\le8$:
\begin{itemize}
\item affine independence of the parent plus all one-flip children;
\item exact-layer cardinalities and constant Hamming weights;
\item Betti numbers of the cumulative simplex skeletons;
\item the relative/mapping-cone rank formula of Theorem~\ref{thm:skeletoncone};
\item the three-variable $G,D,E,F$ calculation extracted from the attached logical-complex notes.
\end{itemize}
The earlier audit scripts separately verify the non-congruence counterexample, the gated Boolean derivative identity, the characteristic-two incidence differential, and the failure of the thesis's printed gluing identity.

\subsection{Carrier-splitting and interaction verification}
The transferred matching gate records 3,242 integral matching cases through
four coarse states, 1,707 product-link checks over $\mathbb F_2$, and 639
additional nonempty matching five-state checks over $\mathbb F_2$.
The interaction gate enumerates 90,201 naturally labelled poset/profile
cases through five states over $\mathbb F_2$, plus 5,000 seeded six-state
trials.  Coarse-poset isomorphism postprocessing yields 87 poset types and
16,761 profiles; exclusive graphs give 70 colour-preserving types with
isolated exclusive vertices retained.  These are different counts of the
same bounded exhaustive evidence, not additional independent cases.

The main interaction run makes 435 selected integral comparisons, supplemented
by nine exact local models, signed integer certificates for both $d^2$
witnesses, and a separate 32-state torsion construction.  The latter checks
181 available links, 31 coarse edge links, and the proved reduced interaction
matrix by Smith normal form, not the whole doubled pair.  The projective-plane
triangulation has facets
\[
012,013,024,035,045,125,134,145,234,235.
\]
An additional 20,000 targeted $K_{2,2}$ trials locate the sharp secondary
differential; these trials are not an exhaustive or uniform isomorphism sample.
All coefficients, seeds, scripts, and integer certificates are recorded in
the accompanying interaction-gate reports.  Computations support the proofs
and do not establish global minimality or publication priority.

\begin{thebibliography}{99}
\small

\bibitem{MatteBlanco1975}
I.~Matte-Blanco,
\emph{The Unconscious as Infinite Sets: An Essay in Bi-Logic},
Duckworth, 1975.

\bibitem{Droncheff2015}
B.~Droncheff,
\emph{On the Ultrametric Unconscious, Quantum Consciousness, and Mental Spaces},
Master's thesis, 2015.

\bibitem{Khrennikov2007}
A.~Yu.~Khrennikov,
``Modelling of psychological behavior on the basis of ultrametric mental space: Encoding of categories by balls,''
\emph{BioSystems} 90 (2007), 656--675.

\bibitem{Murtagh2014}
F.~Murtagh,
``Matte Blanco's bi-logic, metric and ultrametric modelling,''
\emph{Language and Psychoanalysis} 3 (2014).

\bibitem{Dowker1952}
C.~H.~Dowker,
``Homology groups of relations,''
\emph{Annals of Mathematics} 56 (1952), 84--95.

\bibitem{Bjorner2003}
A.~Bj\"orner,
``Nerves, fibers and homotopy groups,''
\emph{Journal of Combinatorial Theory, Series A} 102(1) (2003), 88--93.
DOI: 10.1016/S0097-3165(03)00015-3.

\bibitem{Hatcher2002}
A.~Hatcher,
\emph{Algebraic Topology},
Cambridge University Press, 2002.

\bibitem{Weibel1994}
C.~A.~Weibel,
\emph{An Introduction to Homological Algebra},
Cambridge University Press, 1994.

\bibitem{Curry2013}
J.~Curry,
\emph{Sheaves, Cosheaves and Applications},
Ph.D. thesis, University of Pennsylvania, 2014; arXiv:1303.3255.

\bibitem{Curry2018}
J.~M.~Curry,
``Dualities between cellular sheaves and cosheaves,''
\emph{Journal of Pure and Applied Algebra} 222(4) (2018), 966--993.
DOI: 10.1016/j.jpaa.2017.06.001.

\bibitem{CarlssonDeSilva2010}
G.~Carlsson and V.~de~Silva,
``Zigzag persistence,''
\emph{Foundations of Computational Mathematics} 10 (2010), 367--405.

\bibitem{Edelsbrunner2002}
H.~Edelsbrunner, D.~Letscher, A.~Zomorodian,
``Topological persistence and simplification,''
\emph{Discrete \& Computational Geometry} 28 (2002), 511--533.

\bibitem{CavannaSheehy2018}
N.~J.~Cavanna and D.~R.~Sheehy,
``The generalized persistent nerve theorem,''
arXiv:1807.07920, 2018.

\bibitem{Leitao2026}
A.~Leit\~ao,
``It's All About Covers: Persistent Homology of Cover Refinements,''
arXiv:2602.22784, 2026.

\bibitem{Ramras2026}
D.~A.~Ramras,
``Variations on the Nerve Theorem,''
\emph{Discrete \& Computational Geometry} 75 (2026), 871--903.

\bibitem{GelfandEtAlHypersimplex}
I.~M.~Gelfand, M.~Goresky, R.~MacPherson, V.~Serganova,
``Combinatorial geometries, convex polyhedra, and Schubert cells,''
\emph{Advances in Mathematics} 63 (1987), 301--316.

\bibitem{GrandePadrolSanyal2018}
F.~Grande, A.~Padrol, R.~Sanyal,
``Extension complexity and realization spaces of hyper\-simplices,''
\emph{Discrete \& Computational Geometry} 59 (2018), 621--642.

\bibitem{Gardenfors2000}
P.~G\"ardenfors,
\emph{Conceptual Spaces: The Geometry of Thought},
MIT Press, 2000.


\bibitem{Stong1966}
R.~E.~Stong,
``Finite topological spaces,''
\emph{Transactions of the American Mathematical Society} 123 (1966), 325--340.
DOI: 10.1090/S0002-9947-1966-0195042-2.

\bibitem{Barmak2011}
J.~A.~Barmak,
``On Quillen's Theorem A for posets,''
\emph{Journal of Combinatorial Theory, Series A} 118 (2011), 2445--2453.
DOI: 10.1016/j.jcta.2011.06.008.

\bibitem{BjornerWachsWelker2005}
A.~Bj\"orner, M.~L.~Wachs, and V.~Welker,
``Poset fiber theorems,''
\emph{Transactions of the American Mathematical Society} 357 (2005), 1877--1899.
DOI: 10.1090/S0002-9947-04-03496-8.

\bibitem{FernandezMinian2020}
X.~Fern\'andez and E.~G.~Minian,
``The cylinder of a relation and generalized versions of the Nerve Theorem,''
\emph{Discrete \& Computational Geometry} 63 (2020), 549--559.
DOI: 10.1007/s00454-018-0028-7.

\bibitem{ChacholskiEtAl2021}
W.~Chach\'olski, A.~Jin, M.~Scolamiero, and F.~Tombari,
``Homotopical decompositions of simplicial and Vietoris--Rips complexes,''
\emph{Journal of Applied and Computational Topology} 5 (2021), 215--248.
DOI: 10.1007/s41468-021-00066-2.

\bibitem{DeyHiraniKrishnamoorthy2011}
T.~K.~Dey, A.~N.~Hirani, and B.~Krishnamoorthy,
``Optimal homologous cycles, total unimodularity, and linear programming,''
\emph{SIAM Journal on Computing} 40(4) (2011), 1026--1044.
DOI: 10.1137/100800245.

\bibitem{DroncheffThinning2026}
B.~Theory,
\emph{Binary Order Thinning of Finite Posets: Homology, Collapses, and Sharp Limits},
unpublished companion research manuscript, 27 September 2026, reviewed v3.
The archived title page uses the name Brian Theory.

\bibitem{Bjorner1995}
A.~Bj\"orner,
``Topological methods,''
in \emph{Handbook of Combinatorics}, Vol.~2,
R.~L.~Graham, M.~Gr\"otschel, and L.~Lov\'asz (eds.),
Elsevier, 1995, pp.~1819--1872.

\bibitem{BauerEtAl2023}
U.~Bauer, M.~Kerber, F.~Roll, and A.~Rolle,
``A unified view on the functorial nerve theorem and its variations,''
\emph{Expositiones Mathematicae} 41(4) (2023), 125503.
\href{https://doi.org/10.1016/j.exmath.2023.04.005}{doi:10.1016/j.exmath.2023.04.005}.

\end{thebibliography}

\end{document}
